\section{Project Management and Division of Work}
\label{sec:project-management}

\subsection{Team and roles}

The project was carried out by two MSBT students with Taiko as industry partner. The
roles were separated by artefact rather than by phase, so that each member's contribution
is identifiable in the repository history.

\begin{table}[htbp]
\centering\small
\caption{Division of work.}
\label{tab:roles}
\begin{tabularx}{\textwidth}{@{}l>{\raggedright\arraybackslash}X>{\raggedright\arraybackslash}X@{}}
\toprule
\textbf{Member} & \textbf{Responsibility} & \textbf{Artefacts} \\
\midrule
Yiqun Xu (author) &
Architecture decision to build postage inside TAP; design of all five mechanisms; the
entire implementation, its unit, service and mock-E2E tests; skill and reference
documentation; pull request~\#1; integration of the teammate's evaluation (review and
merge of pull request~\#2). &
25 commits on \path{feat/phase1-notification-coalescing}: 23 implementation commits (\code{1f5f30a}\ldots\code{cee9e0d}; 91 files, $\approx$8{,}200 lines added), one documentation commit (\code{50980f4}) and the merge of pull request~\#2; the \code{@trustedagents/app-postage} package. \\
\addlinespace
Zhang Hanyu (teammate) &
Earlier design exploration (the ``Heed'' track) that preceded the TAP-native decision;
design and execution of the evaluation; documentation of methodology, results and
limitations. &
\path{scripts/benchmark-attention.ts}, \path{scripts/benchmark-postage.ts},
\path{evaluation/results.md}; pull request~\#2 (commit \code{0d123a1}, merged as
\code{0330091}). \\
\bottomrule
\end{tabularx}
\end{table}

This is an individual report: the design rationale, implementation account and
interpretation are the author's; the measurements in Section~\ref{sec:evaluation} are
the teammate's and are reproduced unchanged.

\subsection{Timeline}

\begin{table}[htbp]
\centering\small
\caption{Project timeline as evidenced by the repository.}
\label{tab:timeline}
\begin{tabularx}{\textwidth}{@{}l>{\raggedright\arraybackslash}X@{}}
\toprule
\textbf{When} & \textbf{Milestone} \\
\midrule
Mar--Apr 2026 & Upstream TAP reaches \code{0.2.0-beta.7}: ERC-8004 identity, XMTP transport, grants, \code{tapd} daemon, OpenClaw and Hermes hosts. The upstream maintainers' research note on the messaging layer (23~March) recommends a two-step credit model for anti-spam. Branch point for this work is \code{e52bb66} (21~April). \\
Progress Journal 1 & Project framed as ``Postage: paid attention''; tentative plan is a standalone Tack-like TypeScript service with x402 on Taiko Alethia. \\
May--Aug 2026 & Design work; decision to build the layer inside TAP at the receiver (\S\ref{sec:design-pivot}); implementation and testing of increments 1--5. \\
27--28 Aug 2026 & Increments 1--5 landed on the branch (23 commits). Pull request~\#1 opened against \code{main}. \\
Late Aug 2026 & Industry mentor becomes unreachable (see below). \\
14 Sep 2026 & README, \code{Agents.md} and skill files aligned with the attention/postage features. \\
15 Sep 2026 & Teammate's evaluation merged into the branch (pull request~\#2). \\
18 Oct 2026 & Final report due. \\
\bottomrule
\end{tabularx}
\end{table}

\subsection{Working method}

Work was organised as small, independently green commits in dependency order, each
accompanied by its tests and, where a command or flag changed, by the corresponding
update to the canonical skill file. This mirrored the upstream project's own conventions
(Biome lint and format, strict TypeScript, both E2E suites updated for behavioural change)
and meant that at no point did the branch contain a mechanism without its measurement
hook or its documentation. Design principles (\S\ref{sec:design-principles}) were written
before code and used to settle disagreements with the codebase as they arose; the two
most consequential---``escalations outrank everything'' and ``never lose a message; fail
open''---each vetoed a simpler implementation (a plain cap-and-drop quota) in favour of the
folding design that shipped.

One documentation-alignment commit (\code{50980f4}) was produced with an AI coding
assistant under the author's direction and review; all design decisions and all
production code are the author's own work.

\subsection{Partner engagement and mentor availability}

Taiko's involvement shaped the problem statement: paid interaction between agents,
building on the Tack and x402 integration that already existed in TAP. The credit model
the postage design follows comes from the upstream TAP maintainers' research note, not
from the partner. The industry mentor, Daniel Wang, has been unreachable since late
August~2026. As a consequence:

\begin{itemize}
  \item the decision to replace the Journal-1 architecture with the TAP-native design was
        taken by the team and has \emph{not} been reviewed or endorsed by the partner;
  \item the final artefact carries no mentor sign-off, and this report does not claim one;
  \item the report attributes no position to the partner beyond the problem statement, and
        otherwise presents the team's own reasoning.
\end{itemize}

The author regards this as a material gap in the project's validation rather than a
formality, and has recorded the open questions that would have been put to the mentor in
Section~\ref{sec:future-work}.

\subsection{Risks and how they were handled}

\begin{description}[style=nextline,leftmargin=1.2em]
  \item[Breaking existing users.] Mitigated by making every mechanism opt-in
        (\code{attention.enforce} off by default), carrying stamps in ignorable metadata,
        and leaving the live E2E suite's assertions unchanged.
  \item[Money-moving code with crash windows.] Mitigated by the ordering rules of
        \S\ref{sec:implementation} (pay first, record before transport, complete journal
        after log) and by idempotency tests for each window.
  \item[Platform dependency in evaluation.] The OWS native module's absence on Windows
        blocked Experiment~3. The teammate documented the block rather than substituting a
        weaker measurement; the mechanism's coverage rests on CI tests instead.
  \item[Loss of mentor contact.] Not mitigable within the team; handled by transparency
        in this report and in the pull request description.
\end{description}
